% Source: https://artofproblemsolving.com/community/c3h2784571_2021_abmc_team_rounds_58__actonboxborough_math_competition
% Stable identity: aops-topic-2784571-q13
She bakes at most one pie each year. After a year with a pie she bakes one the next year with probability $\dfrac12$; after a year without, with probability $\dfrac56$. We find the probability of at most one pie in the next $5$ years (this year she baked).

No pies: $\dfrac12\left(\dfrac16\right)^4=\dfrac1{2592}$.

Exactly one pie, in year $j$:
\begin{itemize}
\item $j=1$: $\dfrac12\cdot\dfrac12\cdot\left(\dfrac16\right)^3=\dfrac3{2592}$;
\item $j=2,3,4$: one transition into the pie year ($\dfrac56$), one out of it ($\dfrac12$), the first year without a pie ($\dfrac12$) and two further no-pie years ($\dfrac16$ each), giving $\dfrac12\cdot\dfrac56\cdot\dfrac12\cdot\dfrac1{36}=\dfrac{15}{2592}$ each;
\item $j=5$: $\dfrac12\left(\dfrac16\right)^3\dfrac56=\dfrac5{2592}$.
\end{itemize}
In total $\dfrac{1+3+45+5}{2592}=\dfrac{54}{2592}=\dfrac1{48}$, so the desired probability is $\dfrac{47}{48}$ and $p+q=\boxed{95}$.
